% !TEX root = ../main.tex
\chapter{Topological spaces - construction and purpose}\label{ch:l04}
\begin{paracol}{2}

\begin{sources}
\begin{tabularx}{\linewidth}{@{}l l L@{}}
\toprule
\textbf{Video} & \textbf{Length} & \textbf{Content}\\
\midrule
\seg{4}{1}{Lecture 4: Topological spaces - construction and purpose} & 1:38:42 & Definition of a topology and of a
topological space; how many topologies a finite set has; chaotic, discrete, a three-point and the standard topology on
$\R^d$ (open balls, proof that it is a topology); induced (subset) topology, closed sets, open versus closed, the circle
as a subset and as a quotient; product topology; Furstenberg's topological proof that there are infinitely many primes
(remark); convergence of sequences; a digression on non-transitive dice and ``progress'' in research proposals;
continuity, with examples; homeomorphisms, homeomorphic spaces, and the open classification problem.\\
\bottomrule
\end{tabularx}

\smallskip
\textbf{Transcript} \file{materials/transcripts/C04-V01.txt} (YouTube automatic captions).
\end{sources}

\begin{reading}[Where this lecture sits in the other notes]
\emph{Gravity and Light}, Lecture~1 (Topology) and Lecture~2 (Topological manifolds); \emph{Theoretische Mechanik},
Lecture~2 (Topological manifolds). In both courses topology is the first step after the introduction, and sets,
maps and equivalence relations are taken for granted; here they come from \cref{ch:l02,ch:l03}. (From the lecture
titles; those chapters are not written yet, so the overlap is not checked.)
\end{reading}

% ==================================================================
\section{Topologies and topological spaces}\label{sec:to-def}

Today the lecture starts looking closely at topological spaces, on the basis of the set theory of the previous
lectures \ts{4}{1}{0:04}. As announced last time, establishing a so-called \emph{topology} on a set is the weakest
structure one can put on it in order to have two very important notions: \emph{convergence} of sequences of points of
a set, and \emph{continuity} of maps between two topological spaces \ts{4}{1}{0:25}\ts{4}{1}{0:46}. The definition is at
first sight rather abstract, which is its downside; its upside is that it is extremely simple. It is the result of a
historical development (commented on occasionally), and is to be taken as the simplest definition that mathematicians
in the end found useful \ts{4}{1}{0:46}\ts{4}{1}{1:07}.

\begin{definition}[New concepts: topology, open set]\label{def:to-topology}
Let $M$ be a set (we know what that is from axiomatic set theory). A \textbf{choice} $\mathcal O\subseteq\mathcal P(M)$
of subsets of the power set of $M$ (possibly the entire power set) is called a \textbf{topology} on $M$ if three
conditions hold \ts{4}{1}{1:07}\ts{4}{1}{1:29}\ts{4}{1}{1:50}\ts{4}{1}{2:11}:
\begin{enumerate}
  \item $\emptyset\in\mathcal O$ and $M\in\mathcal O$ \ts{4}{1}{2:31};
  \item $U,V\in\mathcal O\ \Rightarrow\ U\cap V\in\mathcal O$ \ts{4}{1}{2:53}\ts{4}{1}{3:15}\ts{4}{1}{3:35};
  \item for every collection $C\subseteq\mathcal O$ of open sets (arbitrary, not only finite), $\bigcup C\in\mathcal O$
        \ts{4}{1}{4:15}\ts{4}{1}{4:36}.
\end{enumerate}
The elements of $\mathcal O$ are called \textbf{open sets} \ts{4}{1}{21:05}. (The word ``open'' is just a way of speaking
that makes things simple; it gets its meaning in \cref{sec:to-standard}.)
\end{definition}
That is all \ts{4}{1}{5:17}. You are not entirely free in the choice: that is the idea of the three conditions
\ts{4}{1}{2:11}. Condition~2 requires that if you choose $U$ and $V$ you also choose their intersection
\ts{4}{1}{3:35}. The intersection is written with the notation of \cref{ex:se-intersection}: one intersects the elements
of a set, so one can intersect infinitely many sets but never more sets than fit into a set \ts{4}{1}{3:15}\ts{4}{1}{3:55}\ts{4}{1}{4:15}
(``in kindergarten you wrote $U\cap V$'', which is still fine for two sets). Condition~3 could be written in the same
parallel fashion, $\{U,V\}\subseteq\mathcal O\Rightarrow\bigcup\{U,V\}\in\mathcal O$, if one puts the two sets
into a set \ts{4}{1}{4:56}.
\begin{coursenote}[Finite versus arbitrary]
Note the asymmetry of the conditions: \emph{finite} intersections (it suffices to ask for two, then induct), but
\emph{arbitrary} unions. Not said in class.
\end{coursenote}

\begin{definition}[New concepts: topological space]\label{def:to-space}
The pair $(M,\mathcal O)$ of a set and a topology on it is called a \textbf{topological space} \ts{4}{1}{11:17}\ts{4}{1}{11:38}.
\end{definition}
This is the first real instance of what was described very abstractly last time as a \emph{space}: a set with some
additional information, here a choice of subsets \ts{4}{1}{11:38}\ts{4}{1}{12:01}.

\begin{remark}[A real choice: how many topologies are there?]\label{rem:to-count}
Unless $M$ has only one element, there are many different topologies on one and the same set \ts{4}{1}{5:38}\ts{4}{1}{5:59}.
For finite $M$, with $|M|$ the number of elements, the number of topologies is \ts{4}{1}{6:20}\ts{4}{1}{6:42}\ts{4}{1}{7:03}\ts{4}{1}{7:24}:
\begin{center}
\begin{tabular}{l|ccccccc}
$|M|$ & 1 & 2 & 3 & 4 & 5 & 6 & 7\\\hline
topologies & 1 & 4 & 29 & 355 & 6\,942 & 209\,527 & 9\,535\,241
\end{tabular}
\end{center}
For $|M|=1$ the only choice is $\{\emptyset,M\}$. Already for a seven-element set the number is extreme, and an infinite
set has infinitely many topologies \ts{4}{1}{7:45}. So it is a real choice, and depending on it the notions of
continuity of maps and convergence of sequences, which will be defined solely in terms of the topology, change
\ts{4}{1}{8:06}.
\end{remark}
\begin{coursenote}[The numbers]
The captions give the last two numbers as ``2 95 2 7'' and ``9,535,000 241''; the counts of topologies on an
$n$-element set (labelled points) are $1,4,29,355,6942,209527,9535241$, as in the table. The values for 6 and 7 are
as in the standard integer-sequence tables (A000798, checked), not recomputed here.
\end{coursenote}

\subsection*{Examples}

To get familiar with the concept, first the basic examples \ts{4}{1}{8:26}.

\begin{example}[The chaotic topology]\label{ex:to-chaotic}
On any set $M$ take $\mathcal O=\{\emptyset,M\}$ \ts{4}{1}{8:26}. Condition~1 holds trivially. Condition~2: the
intersection of any two elements is $\emptyset\cap\emptyset=\emptyset$, $\emptyset\cap M=\emptyset$, $M\cap M=M$, all
already in $\mathcal O$. Condition~3: the only possible unions are $\emptyset$ or $M$, again in $\mathcal O$
\ts{4}{1}{8:47}\ts{4}{1}{9:08}. It is called the \textbf{chaotic topology} \ts{4}{1}{9:29}. It will not be very useful,
only as an extreme example, ``but sometimes the extreme examples teach you a lot'' \ts{4}{1}{9:52}.
\end{example}

\begin{example}[The discrete topology]\label{ex:to-discrete}
On any set $M$ take $\mathcal O=\mathcal P(M)$, every subset \ts{4}{1}{9:52}\ts{4}{1}{10:14}. The empty set and $M$ are
in the power set; the intersection or the union of any collection of subsets is a subset \ts{4}{1}{10:14}\ts{4}{1}{10:34}.
This is the \textbf{discrete topology}: ``discrete'' in the sense of separated, not in the moral sense
\ts{4}{1}{10:56}.
\end{example}

\begin{example}[A three-point set]\label{ex:to-three}
Let $M=\{1,2,3\}$, which has 29 topologies; choose
$\mathcal O=\bigl\{\emptyset,\{1\},\{2\},\{1,2\},\{1,2,3\}\bigr\}$ \ts{4}{1}{12:21}\ts{4}{1}{12:42}. (The one-element set $\{3\}$ and
the sets containing $3$ other than $M$ are missing.) Condition~1 holds as $\emptyset$ and $M$ are in. Any intersection of
two elements, and any union of elements, again lies in $\mathcal O$: you can check it \ts{4}{1}{13:04}\ts{4}{1}{13:24}.
It has no name, ``not that interesting'', but it is a topology; the chaotic and the discrete ones are two others on the same
set \ts{4}{1}{13:24}\ts{4}{1}{13:46}.
\end{example}

\subsection{The standard topology on $\R^d$}\label{sec:to-standard}

The most heavily used example, which underlies the usual constructions of convergence and continuity on the real
numbers \ts{4}{1}{13:46}\ts{4}{1}{14:06}. Take $M=\R^d=\R\times\dots\times\R$ ($d$ factors; the Cartesian product is known
from set theory) \ts{4}{1}{14:06}\ts{4}{1}{14:27}\ts{4}{1}{14:48}. One could put the chaotic or the discrete topology on it, but there is
one very useful one, the \textbf{standard topology} $\mathcal O_{\rm standard}$ \ts{4}{1}{14:48}\ts{4}{1}{15:09}. In contrast to the
examples above, where the open sets were listed, here they are defined implicitly, as there are uncountably many: the
construction proceeds in three steps \ts{4}{1}{15:09}\ts{4}{1}{15:30}\ts{4}{1}{15:50}.

\paragraph{Step 1: the balls.} For every $x\in\R^d$ and every $r\in\R^+$ define
\[
  B_r(x):=\Bigl\{y\in\R^d\ \Bigm|\ \sqrt{\textstyle\sum_{i=1}^d (y_i-x_i)^2}<r\Bigr\}
\]
(restricted comprehension; $y=(y_1,\dots,y_d)$ is an ordered tuple) \ts{4}{1}{15:50}\ts{4}{1}{16:12}\ts{4}{1}{16:33}\ts{4}{1}{16:54}\ts{4}{1}{17:15}.
Usually this is written as $\|y-x\|<r$, the $2$-norm: squares of the differences, then the square root
\ts{4}{1}{17:36}. At this point it is just a set, nothing to do with topology yet \ts{4}{1}{19:19}; intuitively it is the
\emph{open ball} of radius $r$ around $x$ \ts{4}{1}{18:59}.

One could use the $2n$-norm instead, with the power $2n$ and the $2n$-th root, for some $n$; for the standard topology it
makes no difference. In fact for \emph{any} norm the resulting topology is the same: a famous theorem
\ts{4}{1}{17:56}\ts{4}{1}{18:17}\ts{4}{1}{18:37}. Not to be too general, take $n=1$ \ts{4}{1}{18:37}.

\paragraph{Step 2: the open sets.} A subset $U\subseteq\R^d$ lies in $\mathcal O_{\rm standard}$ by definition if
\[
  \forall p\in U\colon\ \exists r\in\R^+\colon\ B_r(p)\subseteq U
\]
\ts{4}{1}{19:19}\ts{4}{1}{19:40}\ts{4}{1}{20:02}. That is the definition. Remember that ``open'' is just the name of the elements
of a topology \ts{4}{1}{21:05}.

\begin{example}[The picture in $\R^2$]\label{ex:to-picture}
Let the blackboard plane be $\R^2$ and $U$ the interior of a dotted curve, the points on the curve not belonging to $U$
\ts{4}{1}{20:23}\ts{4}{1}{20:44}. Is $U$ open? Pick a point of $U$ blindly; there is a radius such that the ball
around it, drawn dotted too, lies inside $U$ \ts{4}{1}{21:05}\ts{4}{1}{21:25}. The ball consists of all points with distance
\emph{strictly} less than $r$: the boundary does not belong to it \ts{4}{1}{21:45}. Even a point very close to the
boundary has some (small) ball inside $U$ \ts{4}{1}{22:06}. If instead part of the boundary belongs to $U$, take $p$ on
it: however small the positive $r$, part of the ball lies outside $U$, so $U$ is not open \ts{4}{1}{22:27}\ts{4}{1}{22:48}.
\end{example}

\begin{intuition}[Other norms: TV screens and boxes]
With the $2n$-norm for $n=2$ (fourth power, fourth root; one can check it is a norm) the open ``balls'' look a little
like the screens of old TV sets: not round, not square, rounded corners \ts{4}{1}{23:08}\ts{4}{1}{23:29}. They give the same
topology: around a point with a TV set around it one finds a smaller disc, around which a smaller TV set, and so on, so
that what is open for the discs is open for the TV sets and the other way around \ts{4}{1}{23:49}\ts{4}{1}{24:10}. For
$n=10\,000$ the corners are only a little rounded \ts{4}{1}{24:10}\ts{4}{1}{24:31}; for $n\to\infty$ one arrives at
boxes, the \textbf{maximum norm} (``the infinity norm''), writable with absolute values and no infinity
\ts{4}{1}{24:31}\ts{4}{1}{24:52}\ts{4}{1}{25:12}. Every box contains a disc and every disc contains a box, so the topology is always the same
\ts{4}{1}{24:52}\ts{4}{1}{25:12}.
\end{intuition}

\paragraph{Step 3: it is a topology.} It remains to prove that this construction really yields a topology
\ts{4}{1}{25:12}\ts{4}{1}{25:33}.
\begin{proposition}\label{prop:to-standard}
$\mathcal O_{\rm standard}$ on $\R^d$ is a topology.
\end{proposition}
\begin{proof}
Condition~1 \ts{4}{1}{25:57}.
\begin{enumerate}
  \item $\emptyset\in\mathcal O_{\rm standard}$: we must check $\forall p\in\emptyset\colon(\dots)$, with the rest of the
        condition irrelevant. The statement is $\forall p\colon (p\in\emptyset\Rightarrow\dots)$ and $p\in\emptyset$ is always
        false, so the implication is true by ex falso quodlibet, whatever follows; the predicate in $p$ is true for every
        $p$, so the $\forall$-proposition is true \ts{4}{1}{25:57}\ts{4}{1}{26:17}\ts{4}{1}{26:37}\ts{4}{1}{26:58}. The empty set is open
        ultimately because of the definition of the implication.
  \item $\R^d\in\mathcal O_{\rm standard}$: any ball of any radius around any point of $\R^d$ is by definition a subset of
        $\R^d$ \ts{4}{1}{27:18}\ts{4}{1}{27:39}.
\end{enumerate}
Condition~2 \ts{4}{1}{27:39}\ts{4}{1}{28:00}. Let $U,V$ be open; take an arbitrary $p\in U\cap V$ \ts{4}{1}{28:00}\ts{4}{1}{28:20}. Then $p\in U$
and $p\in V$. As $U$ is open there is $r\in\R^+$ with $B_r(p)\subseteq U$; as $V$ is open there is $s\in\R^+$ with
$B_s(p)\subseteq V$ \ts{4}{1}{28:41}\ts{4}{1}{29:01}. Then $B_{\min(r,s)}(p)$ (the smaller of the two, or either one if equal)
lies in both $U$ and $V$, hence in $U\cap V$ \ts{4}{1}{29:23}\ts{4}{1}{29:44}. We found a ball in $U\cap V$ around an arbitrary
point of it, so $U\cap V\in\mathcal O_{\rm standard}$ \ts{4}{1}{30:04}.

Condition~3 \ts{4}{1}{30:04}\ts{4}{1}{30:25}. (``We are not always going to go through all steps; just to demonstrate this once.'')
Let $C\subseteq\mathcal O_{\rm standard}$ be a collection of open sets and $p\in\bigcup C$. Then $p\in U$ for some $U\in C$,
and as $U$ is open there is $r\in\R^+$ with $B_r(p)\subseteq U$. But $U\subseteq\bigcup C$, so $B_r(p)\subseteq\bigcup C$
\ts{4}{1}{30:25}\ts{4}{1}{30:47}\ts{4}{1}{31:08}\ts{4}{1}{31:28}. Hence $\bigcup C$ is open \ts{4}{1}{31:28}\ts{4}{1}{31:49}.
\end{proof}
\begin{recall}[Ex falso quodlibet]
That an implication with false premise is true, so a $\forall$ over an empty range holds, is the argument of the
formal proof that the empty set is unique, \cref{thm:se-empty-unique}.
\end{recall}

% ==================================================================
\section{Constructing new topologies from given ones}\label{sec:to-construct}

Once there is a toolbox of topologies, new ones are mostly custom designed for one's problem and built from the
building blocks one has; it is an important technique to know how to construct new topologies from given ones
\ts{4}{1}{32:14}\ts{4}{1}{32:38}\ts{4}{1}{32:59}.

\subsection{The induced topology on a subset}

\begin{definition}[New concepts: induced (subset) topology]\label{def:to-induced}
Let $(M,\mathcal O)$ be a topological space and $N\subseteq M$ a (proper) subset. Define
\[
  \mathcal O|_N:=\{\,U\cap N\mid U\in\mathcal O\,\}
\]
\ts{4}{1}{32:59}\ts{4}{1}{33:19}\ts{4}{1}{33:40}. Then $\mathcal O|_N$ is a topology on $N$, called the \textbf{induced}
(or \textbf{subset}) \textbf{topology}: ``a definition slash theorem'', since the claim that it is a topology is part of
it \ts{4}{1}{33:40}\ts{4}{1}{34:00}\ts{4}{1}{34:23}.
\end{definition}
It is constructed from the topology of the bigger set by intersecting each open set of it with the subset; every element
is a subset of $N$, so $\mathcal O|_N\subseteq\mathcal P(N)$ and there is a chance that it is a topology
\ts{4}{1}{34:23}\ts{4}{1}{34:45}.

\begin{proof}
It can be boring to go through such proofs, but for the first two or three definitions it is useful for a feeling of how
this works; one always asks the same questions \ts{4}{1}{35:07}.
\begin{enumerate}
  \item $\emptyset\in\mathcal O|_N$: $\emptyset=\emptyset\cap N$ and $\emptyset\in\mathcal O$, so it has the required form
        \ts{4}{1}{35:07}\ts{4}{1}{35:29}.
  \item $N\in\mathcal O|_N$: $N=M\cap N$ and $M\in\mathcal O$ by condition~1 \ts{4}{1}{35:49}\ts{4}{1}{36:10}.
  \item Let $S,T\in\mathcal O|_N$. (Schuller first tried to write it shorter, ran into trouble because he had not used the
        definition first, and repeated more carefully \ts{4}{1}{36:30}\ts{4}{1}{37:15}\ts{4}{1}{37:35}\ts{4}{1}{37:56}.)
        By definition there are $U,V\in\mathcal O$ with $S=U\cap N$ and $T=V\cap N$
        \ts{4}{1}{38:16}\ts{4}{1}{38:36}\ts{4}{1}{38:57}. Then
        \[
          S\cap T=(U\cap N)\cap(V\cap N)=(U\cap V)\cap N,
        \]
        and $U\cap V\in\mathcal O$ as $\mathcal O$ is a topology, so $S\cap T\in\mathcal O|_N$
        \ts{4}{1}{39:18}\ts{4}{1}{39:38}\ts{4}{1}{40:00}.
  \item Arbitrary unions: checked in the same way, ``I am not going to write this down now'' \ts{4}{1}{41:01}
        (compare \cref{ex:to-induced-union}).
\end{enumerate}
\end{proof}
A student asks whether the step $(U\cap N)\cap(V\cap N)=(U\cap V)\cap N$ has been shown: no, it was left for the problem
sheet; it is clear from ``kindergarten set theory'', and with the definition of the intersection from the last problem
sheet the finite case reduces to what one learns there \ts{4}{1}{40:00}\ts{4}{1}{40:20}.
\begin{aside}[Set theory in school]
A digression in class: in the 1970s set theory was a fashion in elementary schools, and parents got unnerved because
``it's mathematics, I don't understand'': red cars, blue cars, and the intersection of red cars and blue cars being empty
\ts{4}{1}{40:20}\ts{4}{1}{40:40}.
\end{aside}
\begin{exercise}[Induced topology]\label{ex:to-induced-union}
Prove the identity $(U\cap N)\cap(V\cap N)=(U\cap V)\cap N$ from the definition of the intersection, and check
condition~3 for $\mathcal O|_N$ (arbitrary unions) \ts{4}{1}{41:01}.
\end{exercise}

\begin{example}[Open in the subset, not open in the big space]\label{ex:to-interval}
Equip $\R$ with the standard topology and take $N=[-1,1]=\{x\in\R\mid -1\le x\le 1\}$ with the induced topology
\ts{4}{1}{41:01}\ts{4}{1}{41:21}\ts{4}{1}{41:43}\ts{4}{1}{42:04}. The set $(0,1]$ is a subset of $\R$; it is \emph{not} open in the
standard topology, because of the point $1$ on its boundary (the argument above) \ts{4}{1}{42:04}\ts{4}{1}{42:26}. But it is
also a subset of $N$, and in the induced topology it \emph{is} open, as
\[
  (0,1]=(0,2)\cap[-1,1],
\]
with $(0,2)$ open in $\R$ \ts{4}{1}{42:26}\ts{4}{1}{42:48}\ts{4}{1}{43:10}\ts{4}{1}{43:31}.
\end{example}
Any set that can be written as an open set of the bigger space intersected with the subspace is open in the induced
topology: a set can be non-open in the big space and nevertheless open in the induced topology of a subset it happens
to lie in \ts{4}{1}{43:31}\ts{4}{1}{43:53}. So the question is always ``open with respect to which topology?'' and one must not
draw quick conclusions from seeing subsets \ts{4}{1}{43:53}\ts{4}{1}{44:14}.

\subsection{Closed sets}

Since ``non-open'' is tempting to call ``closed'', one should say that this is \emph{not} the convention
\ts{4}{1}{44:14}\ts{4}{1}{44:35}.
\begin{definition}[New concepts: closed set]\label{def:to-closed}
Let $(M,\mathcal O)$ be a topological space. A subset $C\subseteq M$ is called \textbf{closed} if its complement
$M\setminus C$ is open \ts{4}{1}{44:35}\ts{4}{1}{44:58}\ts{4}{1}{45:21}.
\end{definition}
\begin{example}[A closed interval]
$[0,1]$ is closed in $\R$ with the standard topology. It has its end points, which only means it is not open; but it is
closed because
\[
  \R\setminus[0,1]=(-\infty,0)\cup(1,\infty)
\]
is a union of two open sets (around each point of either a small ball lies in it), hence open
\ts{4}{1}{45:21}\ts{4}{1}{45:43}\ts{4}{1}{46:04}\ts{4}{1}{46:24}.
\end{example}
\begin{coursenote}[Board notation]
The transcript has ``the interval 0 1'' without brackets; the closed interval $[0,1]$ is meant, as the complement has
open ends at $0$ and $1$ \ts{4}{1}{45:21}\ts{4}{1}{46:04}.
\end{coursenote}

A warning \ts{4}{1}{46:24}\ts{4}{1}{46:46}. For a generic subset of a generic topological space, with nothing known about the
space, six statements are possible: it can be (a) open, (b) closed, (c) open and closed, (d) open and not closed, (e)
not open and closed, (f) not open and not closed \ts{4}{1}{46:46}\ts{4}{1}{47:07}\ts{4}{1}{47:27}\ts{4}{1}{47:48}. Being open and being closed have
nothing to do with each other, except via the complement; from openness one cannot conclude anything about closedness, nor
the other way around \ts{4}{1}{48:09}.
\begin{coursenote}[The six cases]
As listed, (a) and (b) are ``can be open'' and ``can be closed''; the four mutually exclusive cases are
(c) to (f). Schuller's point is that all four combinations of ``open'' and ``closed'' can occur.
\end{coursenote}

\begin{proposition}[Observation: $\emptyset$ and $M$ are open and closed]\label{prop:to-clopen}
In every topological space $(M,\mathcal O)$ the sets $\emptyset$ and $M$ are both open and closed
\ts{4}{1}{48:09}\ts{4}{1}{48:32}.
\end{proposition}
\begin{proof}
$\emptyset=M\setminus M$ and $M\in\mathcal O$, so $\emptyset$ is closed. $M=M\setminus\emptyset$ and
$\emptyset\in\mathcal O$, so $M$ is closed. Both are open by condition~1 \ts{4}{1}{48:32}\ts{4}{1}{48:53}\ts{4}{1}{49:14}.
\end{proof}
There is no contradiction \ts{4}{1}{49:14}\ts{4}{1}{49:36}. Later a \emph{connected} topological space will be defined; one can prove
that in such a space these are the \emph{only} sets that are both open and closed. If one finds this funny, one can look at
connected spaces; but no topology avoids the fact that $\emptyset$ and $M$ are both, as it is a general property
\ts{4}{1}{49:36}\ts{4}{1}{49:56}\ts{4}{1}{50:16}.

\subsection{An application: a topology on the circle}

A real-life example where one uses the induced topology \ts{4}{1}{50:16}. Consider the set $S^1$, the circle, which one
thinks of only as a set: it has no shape; a shape will come only with geometry, at the very end of the building of
structures (``even this is already too much'': it is a \emph{continuous} circle, to come) \ts{4}{1}{50:39}\ts{4}{1}{50:59}\ts{4}{1}{51:19}.

\paragraph{As a subset of $\R^2$.} One can define it as a subset of $\R^2$:
$S^1:=\{(x,y)\in\R^2\mid x^2+y^2=1\}$ \ts{4}{1}{51:19}\ts{4}{1}{51:41}. To get a topology without starting from scratch, let
$\mathcal O_{S^1}:=\mathcal O_{\rm standard}|_{S^1}$ \ts{4}{1}{51:41}\ts{4}{1}{52:04}\ts{4}{1}{52:25}.
\begin{itemize}
  \item An arc of the circle without its two end points is open in this topology: it is $S^1\cap U$ for a suitable open
        $U\in\mathcal O_{\rm standard}(\R^2)$ \ts{4}{1}{52:25}\ts{4}{1}{52:46}\ts{4}{1}{53:07}.
  \item The empty set is open, as always: $\emptyset=\emptyset\cap S^1$ \ts{4}{1}{53:27}.
  \item A set consisting of one point is \emph{not} open: it cannot be written as the circle intersected with an open set of
        $\R^2$. An open set that merely touches the circle has no ball around the point (the boundary is not in it) and
        generates nothing; one that covers the point covers more than one point of the circle
        \ts{4}{1}{53:48}\ts{4}{1}{54:09}\ts{4}{1}{54:29}.
\end{itemize}

\paragraph{As a quotient set.} A different route \ts{4}{1}{54:29}\ts{4}{1}{54:49}: take the real line and identify points under
the equivalence relation $x\sim y:\Leftrightarrow y=x+2\pi$ [?: as heard, ``$y=x+2\pi$''; for $\sim$ to be an equivalence
relation it must be $y=x+2\pi n$ for some $n\in\mathbb Z$, which is what is meant] \ts{4}{1}{54:49}\ts{4}{1}{55:11}. Two points whose
difference is $2\pi$ are identified: starting at $0$, the point $2\pi$ is identified with $0$ and one starts again; this way
the circle is obtained as a quotient set (\cref{def:cl-quotient}) \ts{4}{1}{55:11}\ts{4}{1}{55:31}. One can then study how a topology
is inherited by a quotient set from the topology of the original set; the construction was not shown in class
\ts{4}{1}{55:31}\ts{4}{1}{55:51}.
\begin{exercise}[Quotient topology]\label{ex:to-quotient}
Problem sheet: the topology inherited by a quotient set \ts{4}{1}{55:51}\ts{4}{1}{56:13}. The full statement is in
\cref{ex:to-quotient-full}, \cref{exam:ch04}.
\end{exercise}
So there are sometimes several ways to do the same, and it is convenient to be able to build subsets, quotient sets and the
topologies they inherit \ts{4}{1}{56:13}.

\subsection{The product topology}

One more very important construction \ts{4}{1}{56:13}\ts{4}{1}{56:35}. Let $(A,\mathcal O_A)$ and $(B,\mathcal O_B)$ be topological
spaces (not the induced topology, so no stroke). Very often, and implicitly, one equips the Cartesian product $A\times B$ with
the \textbf{product topology} $\mathcal O_{A\times B}$ \ts{4}{1}{56:35}\ts{4}{1}{56:55}\ts{4}{1}{57:17}.

\begin{definition}[New concepts: product topology]\label{def:to-product}
A set $U\subseteq A\times B$ is an element of $\mathcal O_{A\times B}$ if and only if for every point $p=(a,b)\in U$
there exist $S\in\mathcal O_A$ with $a\in S$ and $T\in\mathcal O_B$ with $b\in T$ such that
$S\times T\subseteq U$ \ts{4}{1}{57:17}\ts{4}{1}{57:40}\ts{4}{1}{58:01}\ts{4}{1}{58:22}\ts{4}{1}{58:42}.
\end{definition}
Intuitively: represent $A$ and $B$ as the two axes (a big assumption, as topological spaces can look very crazy; there is no
notion of dimension yet) \ts{4}{1}{59:04}\ts{4}{1}{59:24}. To check whether a set $U$ is open, take a point $(a,b)$ and look for
an open set $S$ around $a$ in $A$ and $T$ around $b$ in $B$; the box $S\times T$ must lie entirely inside $U$. If it does not
work, choose smaller $S$ and $T$ around the point; whether it then works depends on the topologies of $A$ and $B$
\ts{4}{1}{59:47}\ts{4}{1}{1:00:08}\ts{4}{1}{1:00:28}\ts{4}{1}{1:00:49}. This is inspired by $A,B$ being the real lines with the standard topology; that is the
picture to remember \ts{4}{1}{1:00:49}\ts{4}{1}{1:01:10}.
\begin{exercise}[Product topology]\label{ex:to-product}
Check that $\mathcal O_{A\times B}$ is a topology on $A\times B$ (problem sheet) \ts{4}{1}{1:01:10}\ts{4}{1}{1:01:31}.
\end{exercise}
\begin{remark}[Finite products; $\R^d$]\label{rem:to-product}
Two remarks \ts{4}{1}{1:01:31}\ts{4}{1}{1:01:53}.
\begin{enumerate}
  \item The same works for any finite product $A_1\times A_2\times\dots\times A_n$.
  \item The standard topology $\mathcal O_{\rm standard}$ on $\R^d$ coincides with the product topology of $d$ copies of
        $\R$, each with the standard topology: the argument is the one with the maximum norm, boxes always contain discs and
        discs always contain boxes \ts{4}{1}{1:01:53}\ts{4}{1}{1:02:15}\ts{4}{1}{1:02:36}.
\end{enumerate}
One could have been more minimalistic, defining the standard topology only on $\R$ and then on $\R^d$ via the product topology.
Both ways coincide; there is no a priori reason why they should, but they do \ts{4}{1}{1:02:36}\ts{4}{1}{1:02:57}\ts{4}{1}{1:03:17}.
\end{remark}

\begin{secondary}[Furstenberg: topology in number theory]
A remark before continuing \ts{4}{1}{1:03:17}\ts{4}{1}{1:03:38}. In this course topologies give structure to sets on the way to
manifolds and physical space. But topology is not only useful in this ``roughly geometric'' context: it is a very versatile
subject and deserves its own status \ts{4}{1}{1:03:38}\ts{4}{1}{1:03:58}. For instance one can equip the integers with a (cleverly
chosen) topology: Furstenberg discovered that with a suitable topology on $\mathbb Z$ one obtains a purely topological proof, by
just thinking about open and closed sets, that there are infinitely many primes \ts{4}{1}{1:03:58}\ts{4}{1}{1:04:20}\ts{4}{1}{1:04:42}. So
even in number theory topology is a tool. The example is famous as it is unexpected, and it explains why the definition has this
abstract form: it covers many situations, and the brain work goes into choosing the right topology to solve one's
problem, which sometimes takes ingenuity \ts{4}{1}{1:04:42}\ts{4}{1}{1:05:03}\ts{4}{1}{1:05:23}.\par
\smallskip
\emph{Not said in class.} Furstenberg's topology on $\mathbb Z$ has the arithmetic progressions as a basis; the proof is a
page long.
\end{secondary}
\begin{history}[Hausdorff's neighbourhoods]
\emph{Not said in class} (the lecture only says that the definition came from a historical development, \ts{4}{1}{0:46}\ts{4}{1}{1:07}).\par
\smallskip
\emph{Checked.} The notion of an abstract topological space with neighbourhood axioms is due to F.~Hausdorff,
\emph{Grundz\"uge der Mengenlehre} (1914); the definition by open sets, as above, became standard later. H.~Furstenberg,
``On the infinitude of primes'', \emph{Amer.\ Math.\ Monthly} 62 (1955) 353.
\end{history}

% ==================================================================
\section{Convergence}\label{sec:to-conv}

The whole structure is only useful if one can do something with it: convergence of sequences and continuity of maps
\ts{4}{1}{1:05:23}\ts{4}{1}{1:05:45}.

\begin{definition}[New concepts: sequence, convergence, open neighbourhood]\label{def:to-conv}
A \textbf{sequence} on a topological space $(M,\mathcal O)$ is a map $q\colon\N\to M$, from the non-negative integers into the
set underlying the space (Schuller first said ``integers'' and corrected to the naturals) \ts{4}{1}{1:06:09}\ts{4}{1}{1:06:33}\ts{4}{1}{1:08:00}. A sequence
$q$ \textbf{converges} against a point $a\in M$, the \textbf{limit point}, if for every open set $U\in\mathcal O$ with $a\in U$
there is $N\in\N$ such that
\[
  \forall n>N\colon\ q(n)\in U
\]
\ts{4}{1}{1:06:33}\ts{4}{1}{1:06:58}\ts{4}{1}{1:07:19}\ts{4}{1}{1:08:00}\ts{4}{1}{1:08:21}.
\end{definition}
An open set $U$ with $a\in U$ is called an \textbf{open neighbourhood} of $a$: it lies around the point (what lies around you is
your neighbourhood; this is just topology lingo) \ts{4}{1}{1:07:19}\ts{4}{1}{1:07:39}. The condition must hold for \emph{every} open neighbourhood,
which is very important \ts{4}{1}{1:08:00}: from some $N$ on, all the images of the sequence lie in it
\ts{4}{1}{1:08:00}\ts{4}{1}{1:08:21}.

\begin{example}[Chaotic topology: everything converges to everything]\label{ex:to-conv-chaotic}
Let $M$ carry the chaotic topology \ts{4}{1}{1:08:43}. Being a sequence has nothing to do with the topology; it is a map into the
set \ts{4}{1}{1:09:05}. Claim: \emph{every sequence converges against every point} \ts{4}{1}{1:09:25}. Indeed let $q$ be arbitrary and
$a\in M$ arbitrary \ts{4}{1}{1:09:48}. Of the two open sets only $M$ contains $a$, so we must check $\exists N\ \forall n>N\colon q(n)\in M$,
which holds for every $n$ as $q$ maps into $M$, independently of the point and of the sequence
\ts{4}{1}{1:10:09}\ts{4}{1}{1:10:29}\ts{4}{1}{1:10:50}.
\end{example}
Hence ``chaotic'', and not useful \ts{4}{1}{1:10:50}.
\begin{aside}[A quick exam solution]
Schuller's joke: if an examiner asks to show that a sequence converges but is not careful enough to specify the topology, you
can claim that it converges against every point, since you choose the chaotic topology and were not asked to choose another one
\ts{4}{1}{1:10:50}\ts{4}{1}{1:11:11}.
\end{aside}

\begin{example}[Discrete topology: only almost constant sequences]\label{ex:to-conv-discrete}
Equip $M$ with the discrete topology \ts{4}{1}{1:11:33}. Every set is open, in particular every $\{a\}$ is an open
neighbourhood of $a$, so a sequence converges against $a$ only if $q(n)=a$ for all large $n$. So only \emph{almost constant}
sequences converge, and they converge against that constant; ``almost constant'' means constant everywhere apart from a finite
exception set \ts{4}{1}{1:11:33}\ts{4}{1}{1:11:53}\ts{4}{1}{1:12:14}.
\end{example}

\begin{example}[$\R^d$ with the standard topology]\label{ex:to-conv-standard}
Let $M=\R^d$ with $\mathcal O_{\rm standard}$ \ts{4}{1}{1:12:37}. One can prove that a sequence $q$ converges against $a\in\R^d$
if and only if
\[
  \forall\varepsilon>0\ \exists N\in\N\ \forall n>N\colon\ \|q(n)-a\|<\varepsilon,
\]
for some (any) norm \ts{4}{1}{1:12:37}\ts{4}{1}{1:13:00}\ts{4}{1}{1:13:21}\ts{4}{1}{1:13:42}.
\end{example}
\begin{exercise}[The $\varepsilon$ criterion]\label{ex:to-conv-eps}
Prove the equivalence in \cref{ex:to-conv-standard} from \cref{def:to-conv} and the definition of $\mathcal O_{\rm standard}$
(stated in class without proof).
\end{exercise}
A (garbled) question from a student, apparently about the relation with the analysis way of saying ``for all\dots'' [?], is answered
as follows \ts{4}{1}{1:13:42}\ts{4}{1}{1:14:03}. The discrete topology is the \emph{finest} one: everything convergent in it is convergent
in any other topology, but in topologies that are not that fine, with fewer subsets chosen, there are more convergent
sequences, which do not converge in the finer one \ts{4}{1}{1:13:42}\ts{4}{1}{1:14:03}\ts{4}{1}{1:14:23}.

\begin{example}[$q(n)=1-1/n$]\label{ex:to-conv-1n}
In $M=\R$ take $q(n)=1-\frac1n$ for $n\ge1$ (start at $n=1$ not to run into trouble at zero) \ts{4}{1}{1:14:23}\ts{4}{1}{1:14:45}. On the real line the
points are $0,\frac12,\frac23,\dots$, closer and closer to $1$ \ts{4}{1}{1:14:45}\ts{4}{1}{1:15:05}. Every value is different, so the sequence is not almost constant, and
it does \emph{not} converge in $\R$ with the discrete topology; but it converges against $1$ in $\R$ with the standard
topology \ts{4}{1}{1:15:05}\ts{4}{1}{1:15:34}\ts{4}{1}{1:15:57}.
\end{example}
The intuition of ``getting closer and closer'' is that of the standard topology, which uses the norm, i.e.\ the distance
\ts{4}{1}{1:15:57}\ts{4}{1}{1:16:18}. It generalises to every \emph{metric space}: from a metric (a distance function) one can construct a
topology; but by far not every topology is induced by a distance map \ts{4}{1}{1:16:18}\ts{4}{1}{1:16:39}.

\begin{secondary}[Digression: research proposals, non-transitive dice and topologies of theories]
A half-serious side remark \ts{4}{1}{1:16:39}\ts{4}{1}{1:17:00}. To get research money one writes a proposal of the type ``we do this
and that, because then we will know more''; in physics: ``we come closer to the truth'', or ``in recent years we have converged
against the true theory''. After the (spiritual) approval of the office for gender and diversity it goes to experts, who have no time
and read only the introduction, which must say that we are getting better and making progress
\ts{4}{1}{1:17:00}\ts{4}{1}{1:17:20}\ts{4}{1}{1:17:40}.

What would this mean? One way to say that theories make progress is that they get better and better; for that one needs a measure
of whether one theory is better than another, a quality function $Q$ \ts{4}{1}{1:17:40}\ts{4}{1}{1:18:01}. Consider \textbf{non-transitive dice}
(Schuller said ``non-commutative'' and corrected himself): three dice $a,b,c$ with different distributions of numbers on their
faces (the distribution on the board was only illustrative) \ts{4}{1}{1:18:01}\ts{4}{1}{1:18:27}\ts{4}{1}{1:18:48}. Two dice play
against each other: both are thrown, the one with the higher number gets a point, and after a long run the one with more points
wins \ts{4}{1}{1:19:08}\ts{4}{1}{1:19:29}\ts{4}{1}{1:19:51}. One can calculate with which probability each wins, so one could say that
the quality of $a$ was greater than that of $b$ if $a$ wins against $b$ \ts{4}{1}{1:19:51}\ts{4}{1}{1:20:12}. Similarly $b$ may win
against $c$. Then one would conclude that $a$ wins against $c$, but remarkably there is a distribution of numbers on the three dice such
that $c$ wins against $a$: from $Q(a)>Q(b)$ and $Q(b)>Q(c)$ one cannot conclude $Q(a)>Q(c)$ \ts{4}{1}{1:20:12}\ts{4}{1}{1:20:34}\ts{4}{1}{1:20:54}.
So no such real-valued quality measure exists \ts{4}{1}{1:20:54}\ts{4}{1}{1:21:14}.

This is remarkable: anybody who assigns to complex situations, such as research proposals or the abilities of a girl or a boy, a
real-valued function measuring them on one scale, and compares on it, makes an assumption on which economics, for instance, is built
\ts{4}{1}{1:21:14}\ts{4}{1}{1:21:35}. Already the little complexity of numbers on dice has no quality function; so nobody should say
``this theory is better than that theory'' without being asked immediately for the quality function, which is generically
non-transitive; only special distributions are transitive \ts{4}{1}{1:21:35}\ts{4}{1}{1:21:57}.

One might say: even without a quality function, the theories come \emph{closer to the true theory}. That would be a very weak but
still impressive statement; but to make it one needs to know a \emph{topology on the space of all theories}, which is far too
complex; Schuller does not know how to write it down, and assumes, unless somebody convinces him otherwise, that statements of the
type ``theories converge somehow'' are not meaningful \ts{4}{1}{1:21:57}\ts{4}{1}{1:22:17}\ts{4}{1}{1:22:38}\ts{4}{1}{1:22:58}. Since no
structure making fewer assumptions than a topology is known by which to define convergence, topology also reveals nonsensical or
ill-defined statements that sound good; it raises at least sufficient suspicion about any such claim, and applies to other
circumstances too \ts{4}{1}{1:22:58}\ts{4}{1}{1:23:20}\ts{4}{1}{1:23:41}.
\end{secondary}
\begin{aside}[Hairdressers]
A passing joke: ``what hairdressers can do, only hairdressers can do; in logic, in topology, maybe you can conclude\dots''
\ts{4}{1}{1:23:20}. (Captions garbled; the point is the one of the digression above.) [?]
\end{aside}
\begin{history}[Efron's dice]
\emph{Not said in class.} \emph{Checked.} Sets of non-transitive dice were found by the statistician Bradley Efron and were
popularised by Martin Gardner in \emph{Scientific American} (Dec.\ 1970).
\end{history}

% ==================================================================
\section{Continuity}\label{sec:to-cont}

Convergence was interesting, but continuity is at the heart of topology, as it will finally be used to define the structure
preserving maps \ts{4}{1}{1:23:41}\ts{4}{1}{1:24:02}.

\begin{definition}[New concepts: continuous map]\label{def:to-cont}
Let $(M,\mathcal O_M)$ and $(N,\mathcal O_N)$ be topological spaces and $f\colon M\to N$ a map between the underlying sets. Then $f$
is called \textbf{continuous} if for every open set $V\in\mathcal O_N$ in the target space the pre-image $\operatorname{preim}_f(V)$
(\cref{def:cl-preimage}) is an open set in the domain, $\operatorname{preim}_f(V)\in\mathcal O_M$
\ts{4}{1}{1:24:02}\ts{4}{1}{1:24:23}\ts{4}{1}{1:24:43}\ts{4}{1}{1:25:07}\ts{4}{1}{1:25:29}\ts{4}{1}{1:25:50}.
\end{definition}
Extremely simple: ``pre-images of open sets are open''. That hides the fact that the first openness is in the target space and the second in the domain, so more precisely
$\forall V\in\mathcal O_N\colon\operatorname{preim}_f(V)\in\mathcal O_M$ \ts{4}{1}{1:25:50}\ts{4}{1}{1:26:14}\ts{4}{1}{1:26:39}. It is at the heart of topology, as will be
seen \ts{4}{1}{1:26:39}.
\begin{recall}[Pre-image]
The pre-image of $V\subseteq N$ is the set of all points $m\in M$ with $f(m)\in V$, \cref{def:cl-preimage}
\ts{4}{1}{1:25:29}.
\end{recall}

\begin{example}[Discrete domain: every map is continuous]\label{ex:to-cont-discrete}
Take any map $f\colon M\to N$ between arbitrary sets, with the discrete topology on $M$ and any topology on $N$ \ts{4}{1}{1:26:39}\ts{4}{1}{1:27:00}.
Then $f$ is continuous \ts{4}{1}{1:27:20}: for any open $V\subseteq N$ the pre-image is a subset of $M$, and every subset of $M$ is open (for $V=\emptyset$ the
pre-image is the empty set; no need to distinguish cases) \ts{4}{1}{1:27:20}\ts{4}{1}{1:27:40}\ts{4}{1}{1:28:01}. No property of $f$ was used.
\end{example}
If every map is continuous, continuity is a useless property; so one \emph{never} equips the domain with the discrete topology
if one still wants to work with continuity \ts{4}{1}{1:28:01}\ts{4}{1}{1:28:21}.

\begin{example}[Chaotic target: every map is continuous]\label{ex:to-cont-chaotic}
Take any map $M\to N$ with any topology on $M$ and the chaotic topology on $N$: again every map is continuous, as one can quickly
think about (the only non-empty open set of $N$ is $N$, with pre-image $M$, and $\emptyset$ has pre-image $\emptyset$)
\ts{4}{1}{1:28:21}\ts{4}{1}{1:28:42}\ts{4}{1}{1:29:02}.
\end{example}

\begin{example}[$\R^d\to\R^f$: the $\varepsilon$--$\delta$ criterion]\label{ex:to-cont-eps}
For $f\colon\R^d\to\R^{f'}$ with both spaces carrying the standard topology one recovers the usual definition of continuity: for
every $\varepsilon>0$ there is $\delta>0$ such that for all $x,y$ with $\|x-y\|<\delta$ one has $\|f(x)-f(y)\|<\varepsilon$ \ts{4}{1}{1:29:02}\ts{4}{1}{1:29:24}. The
$\varepsilon$--$\delta$ criterion is recovered from the general definition, which makes sense for any pair of topological spaces
\ts{4}{1}{1:29:24}\ts{4}{1}{1:29:44}.
\end{example}
\begin{coursenote}[A correction on the board]
A student points at the target dimension (written with the letter of the map): it must be an $n$ (here $f'$ to avoid a clash with the map $f$)
\ts{4}{1}{1:29:44}\ts{4}{1}{1:30:04}. ``Thank you very much, yes of course.''
\end{coursenote}
\begin{exercise}[$\varepsilon$--$\delta$]\label{ex:to-cont-eps-proof}
Derive the $\varepsilon$--$\delta$ criterion from \cref{def:to-cont} for $\mathcal O_{\rm standard}$ on $\R^d$ and $\R^{f'}$ (stated in
class without proof).
\end{exercise}

% ==================================================================
\section{Homeomorphisms: structure preserving maps of topological spaces}\label{sec:to-homeo}

\begin{definition}[New concepts: homeomorphism, homeomorphic]\label{def:to-homeo}
Let $f\colon M\to N$ be a bijection (\cref{def:cl-bij}) between two sets and make topological spaces out of them with $\mathcal O_M$ and
$\mathcal O_N$. Then $f$ is called a \textbf{homeomorphism} if \ts{4}{1}{1:30:04}\ts{4}{1}{1:30:26}\ts{4}{1}{1:30:47}
\begin{enumerate}
  \item $f\colon M\to N$ is continuous, and
  \item the inverse $f^{-1}\colon N\to M$ (it exists and is unique since $f$ is a bijection, \cref{def:cl-inverse}) is continuous
        \ts{4}{1}{1:31:08}\ts{4}{1}{1:31:30}.
\end{enumerate}
Two topological spaces between which a homeomorphism exists are called \textbf{homeomorphic} \ts{4}{1}{1:35:07}, and one writes
$(M,\mathcal O_M)\cong_{\rm top}(N,\mathcal O_N)$ [?: the same isomorphism sign as in set theory is used].
\end{definition}
The word ``homeomorphism'' is long, so one says ``homeo'' and suppresses ``morphism'' \ts{4}{1}{1:31:30}\ts{4}{1}{1:31:51}. The origin of the name:
\emph{homoios} (Greek), similar, and \emph{morphe}, shape (not said in class).

Homeomorphisms are \emph{the structure preserving maps} in topology, the central notion of map \ts{4}{1}{1:30:47}\ts{4}{1}{1:31:51}. Why? $f$ takes
one from $M$ to $N$ as a set, and $f^{-1}$ back; since $f$ is continuous the pre-image of any open set is open, and the
pre-image under $f^{-1}$ is the image under $f$; so, if a homeomorphism exists, it provides a one-to-one pairing of the open sets of $M$ with those of $N$
according to the chosen topologies \ts{4}{1}{1:31:51}\ts{4}{1}{1:32:15}\ts{4}{1}{1:32:36}\ts{4}{1}{1:32:57}\ts{4}{1}{1:33:21}. One can say that the topologies
have the same structure \ts{4}{1}{1:33:21}\ts{4}{1}{1:33:44}.

\begin{proposition}[Homeomorphisms pair up open sets]\label{prop:to-homeo-pair}
Let $f\colon M\to N$ be a homeomorphism. Then $U\mapsto f(U)$ is a bijection $\mathcal O_M\to\mathcal O_N$ with inverse $V\mapsto f^{-1}(V)$.
\end{proposition}
\begin{proof}
(Not written out in class; the content of the argument above.) For $U\in\mathcal O_M$, $f(U)$ is the pre-image of $U$ under the continuous map
$f^{-1}$, hence open in $N$; for $V\in\mathcal O_N$, $f^{-1}(V)$ is the pre-image under the continuous map $f$, so open in $M$. The two assignments are inverse to each other
because $f$ is a bijection.
\end{proof}

\begin{intuition}[Isomorphic as sets, different as topological spaces]\label{int:to-hair}
A bijection exists, so the two sets are isomorphic as sets, ignoring the topology; if they are also homeomorphic they are isomorphic as topological
spaces, which is more \ts{4}{1}{1:33:44}\ts{4}{1}{1:34:05}\ts{4}{1}{1:34:26}\ts{4}{1}{1:34:47}. The same isomorphism sign as in set theory is used, but the bijection
now has to satisfy the additional homeomorphism condition. Conversely, topological isomorphism implies set-theoretic isomorphism; but two sets that are
isomorphic can, once equipped with topologies, become non-isomorphic (different topologically) \ts{4}{1}{1:34:47}\ts{4}{1}{1:35:07}\ts{4}{1}{1:35:30}\ts{4}{1}{1:35:51}.
Things that are the same up to reshuffling can become non-isomorphic at a higher structural level by establishing different structures
\ts{4}{1}{1:35:51}. ``Two children, both blonde; now I dye their hair, one blonde, the other red: now they are different. One finds other distinguishing
properties, but fine'' \ts{4}{1}{1:36:13}.
\end{intuition}

\paragraph{Classification.} Recall the recurrent theme of mathematics, \cref{int:cl-theme}: classify spaces by the structure preserving maps between them
\ts{4}{1}{1:36:13}\ts{4}{1}{1:36:33}. In set theory one classifies by cardinality (not carried far, there is a very big theory behind it); for finite sets, by the number of
elements: two finite sets with the same number of elements are isomorphic as sets \ts{4}{1}{1:36:33}\ts{4}{1}{1:36:53}.
For topological spaces one would like properties that distinguish them in the same way; such properties will be defined in the next lecture
\ts{4}{1}{1:36:53}\ts{4}{1}{1:37:13}. But nobody knows a \emph{complete} set of properties, like the number of elements for sets, such that spaces with the same
values of all of them are already homeomorphic \ts{4}{1}{1:37:13}\ts{4}{1}{1:37:34}\ts{4}{1}{1:37:55}. There is no classification of topological spaces as far as Schuller
knows (``there are just too many, or we haven't thought carefully enough''); this is an open classification problem. One can prove that, ``and then make donations
to the university'': it is important to know what the open problems of a field are, to solve one of them one had better know they are open
\ts{4}{1}{1:37:55}\ts{4}{1}{1:38:15}.
\begin{coursenote}[On the open problem]
Not said in class: the homeomorphism problem is undecidable for manifolds of dimension $\ge4$ (Markov, 1958), so no complete, algorithmic list
of invariants exists there; for general topological spaces the situation is wilder still. Stated here from memory, not checked.
\end{coursenote}

% ==================================================================
\section*{Key concepts}
\begin{keyconcepts}
\textbf{Retrieval questions.}
\begin{enumerate}
  \item State the three conditions of a topology. Which are finite, which arbitrary?
  \item Why are the chaotic, the discrete and the $\{\emptyset,\{1\},\{2\},\{1,2\},\{1,2,3\}\}$ examples topologies?
  \item How is the standard topology on $\R^d$ built in three steps, and why does the choice of norm not matter?
  \item Why is $(0,1]$ open in $[-1,1]$ but not in $\R$? Why is $\{p\}$ not open in $S^1\subseteq\R^2$?
  \item Define closed sets; why are $\emptyset$ and $M$ both open and closed?
  \item Which sequences converge in the chaotic, which in the discrete topology? Does $1-\frac1n$ converge in the standard one?
  \item Define continuity. Why is every map continuous out of a discrete space and into a chaotic one?
  \item Define homeomorphism; why does the classification of topological spaces remain open?
\end{enumerate}
\textbf{Concepts.}
\begin{itemize}
  \item Topology, open sets, topological space (\cref{def:to-topology,def:to-space}); standard topology on $\R^d$ (\cref{prop:to-standard}).
  \item Induced topology $\mathcal O|_N=\{U\cap N\}$ (\cref{def:to-induced}); closed set $=$ open complement (\cref{def:to-closed,prop:to-clopen}).
  \item Product topology (\cref{def:to-product}); circle as subset and as quotient.
  \item Convergence: every open neighbourhood contains the tail (\cref{def:to-conv}).
  \item Continuity: pre-images of open sets are open (\cref{def:to-cont}); homeomorphism (\cref{def:to-homeo}).
\end{itemize}
\end{keyconcepts}

\flushside
\end{paracol}
